\begin{figure}[htbp]
\centering
\begin{tikzpicture}
\begin{axis}[
    width=14cm, height=9cm,
    xlabel={$x$},
    ylabel={Probability Density $f(x)$},
    xmin=-4.5, xmax=4.5,
    ymin=0, ymax=0.5,
    xtick={-4,-3,-2,-1,0,1,2,3,4},
    xticklabels={$\mu-4\sigma$,$\mu-3\sigma$,$\mu-2\sigma$,$\mu-\sigma$,$\mu$,$\mu+\sigma$,$\mu+2\sigma$,$\mu+3\sigma$,$\mu+4\sigma$},
    x tick label style={font=\small},
    grid=major,
    grid style={gray!30},
    axis on top,
    major tick length=0.08cm,
    clip=false
]

% Fill regions for 68-95-99.7 rule (from widest to narrowest)
% 99.7% region (3 sigma)
\addplot[fill=UofSCHorseshoe!25, draw=none, domain=-3:3, samples=100]
    {1/(sqrt(2*pi)) * exp(-x^2/2)} \closedcycle;

% 95% region (2 sigma)
\addplot[fill=UofSCAtlantic!35, draw=none, domain=-2:2, samples=100]
    {1/(sqrt(2*pi)) * exp(-x^2/2)} \closedcycle;

% 68% region (1 sigma)
\addplot[fill=UofSCGarnet!40, draw=none, domain=-1:1, samples=100]
    {1/(sqrt(2*pi)) * exp(-x^2/2)} \closedcycle;

% Main Gaussian curve
\addplot[UofSCBlack, very thick, domain=-4.5:4.5, samples=200]
    {1/(sqrt(2*pi)) * exp(-x^2/2)};

% Vertical lines at sigma boundaries
\draw[UofSCGarnet, thick, dashed] (-1, 0) -- (-1, 0.242);
\draw[UofSCGarnet, thick, dashed] (1, 0) -- (1, 0.242);

\draw[UofSCAtlantic, thick, dashed] (-2, 0) -- (-2, 0.054);
\draw[UofSCAtlantic, thick, dashed] (2, 0) -- (2, 0.054);

\draw[UofSCHorseshoe, thick, dashed] (-3, 0) -- (-3, 0.0044);
\draw[UofSCHorseshoe, thick, dashed] (3, 0) -- (3, 0.0044);

% Mean line
\draw[UofSCBlack, thick] (0, 0) -- (0, 0.42);
\node[UofSCBlack, font=\bfseries, above] at (0, 0.42) {Mean ($\mu$)};

% 68% label
\draw[UofSCGarnet, thick, <->] (-1, 0.32) -- (1, 0.32);
\node[UofSCGarnet, font=\bfseries, above] at (0, 0.32) {68.27\%};

% 95% label
\draw[UofSCAtlantic, thick, <->] (-2, 0.42) -- (2, 0.42);
\node[UofSCAtlantic, font=\bfseries, above] at (0, 0.42) {95.45\%};

% 99.7% label
\draw[UofSCHorseshoe, thick, <->] (-3, 0.48) -- (3, 0.48);
\node[UofSCHorseshoe, font=\bfseries, above] at (0, 0.48) {99.73\%};

% Tail probability annotations
\node[UofSCCongaree, font=\footnotesize, align=center] at (-3.8, 0.12) {Tail\\0.135\%};
\node[UofSCCongaree, font=\footnotesize, align=center] at (3.8, 0.12) {Tail\\0.135\%};

% PDF formula
\node[UofSCBlack, font=\small, align=left] at (-3.3, 0.38) {$f(x) = \frac{1}{\sigma\sqrt{2\pi}}\exp\left(-\frac{(x-\mu)^2}{2\sigma^2}\right)$};

% Standard deviation arrows
\draw[UofSCRose, thick, |<->|] (0, -0.03) -- (1, -0.03);
\node[UofSCRose, font=\footnotesize, below] at (0.5, -0.03) {$\sigma$};

\end{axis}
\end{tikzpicture}
\caption{The Gaussian (normal) distribution illustrating the 68-95-99.7 empirical rule. The probability density function is shown with shaded regions indicating the probability mass within one (Garnet, 68.27\%), two (Atlantic blue, 95.45\%), and three (Horseshoe green, 99.73\%) standard deviations of the mean. This rule is fundamental in statistical quality control, confidence interval estimation, and Kalman filter design. In control systems, sensor noise is often modeled as Gaussian with known $\mu$ and $\sigma$, making these probability regions essential for uncertainty quantification.}
\label{fig:normal_curve}
\end{figure}
